\chapter{Beyond the first phase}
\label{ch:beyond}

The nodes in this chapter are not deliverables of the first phase. They record the road
ahead so that the design choices above are made with them in view; each is marked as not
ready.

\begin{theorem}[Triangulated structure]
\label{thm:derived-triangulated}
\notready
\uses{def:dg-derived-category, prop:resolutions, def:dg-module-shift}
$\Ch(\cA)$ with the split conflations is a Frobenius category whose stable category is
$\Hcat(\cA)$; the suspension is $M \mapsto M[1]$, and $\Dcat(\cA)$ inherits a triangulated
structure in which short exact sequences of dg-modules give triangles (Keller, Lemma 3.3).
Stretch goal of the first phase.
\end{theorem}

\begin{theorem}[Model structures on dg-modules]
\label{thm:model-structures}
\notready
\uses{def:dg-module-cat, def:dg-qis, def:cofibrant-fibrant}
$\Ch(\cA)$ carries a projective and an injective Quillen model structure with weak
equivalences the quasi-isomorphisms (Keller, Theorem 3.2). Both are cofibrantly generated;
this is where the recognition theorem for cofibrantly generated model categories, developed
in the PI's parallel project, enters.
\end{theorem}

\begin{theorem}[Compact generators]
\label{thm:compact-generators}
\notready
\uses{thm:derived-triangulated, cor:derived-homs}
$\Dcat(\cA)$ has arbitrary coproducts and the shifted representables $X^{\wedge}[n]$ form a
set of compact generators; the compact objects are the perfect dg-modules
(Keller, Corollary 3.7).
\end{theorem}

\begin{definition}[Bimodules and tensor functors]
\label{def:bimodules}
\notready
\uses{def:dg-tensor, def:dg-module}
An $\cA$-$\cB$-bimodule is a dg-module over $\cA^{\op} \otimes \cB$; it induces an adjoint
pair between $\Ch(\cA)$ and $\Ch(\cB)$ whose derived functors relate $\Dcat(\cA)$ and
$\Dcat(\cB)$ (Keller \S 3.8).
\end{definition}

\begin{theorem}[dg Morita equivalence]
\label{thm:dg-morita}
\notready
\uses{def:bimodules, thm:compact-generators, def:quasi-equivalence}
Keller's Theorem 3.11 and Rickard's Theorem 3.12: characterizations of when
$\Dcat(\cA) \simeq \Dcat(\cB)$ is induced by tensor functors, and of derived equivalence
for algebras with homology in degree $0$.
\end{theorem}

\begin{theorem}[The homotopy category of dg-categories]
\label{thm:dgcat-model}
\notready
\uses{def:quasi-equivalence, def:dg-tensor, def:dg-functor-category, thm:model-structures}
Tabuada's cofibrantly generated model structure on $\dgcat$ with weak equivalences the
quasi-equivalences, and Toën's description of the mapping spaces of its homotopy category
through bimodules, giving internal Homs in $\mathrm{Ho}(\dgcat)$ (Keller \S 4). This is the
long-range target: derived Morita theory in the unbounded setting.
\end{theorem}
